\section{Results and Discussion}

This section presents the results of the multi-lane intersection simulation, validating the matrix-based model against theoretical predictions and exploring the sensitivity of congestion to key parameters.

\subsection{Model Validation}

The simulation was validated by comparing the computed average delay against the theoretical delay formula (Equation 8):

\begin{equation}
W = \frac{r^2}{2C}\left(1 + \frac{\lambda}{s - \lambda}\right)
\end{equation}

where $r$ is the red time, $C = 2g + b$ is the cycle length (with $g$ as green time per direction and $b$ as bike phase duration), $\lambda$ is the arrival rate, and $s$ is the saturation flow rate calculated from Equations (2)--(4).

\subsubsection{Baseline Parameters}

The web simulator now mirrors the six-stage cycle described in Section~6 of the paper. Unless noted otherwise, the following defaults are used in the experiments below:
\begin{itemize}
    \item Lane configuration: three approach lanes per leg (dedicated left, through, and right turns) with the geometry shown in Fig.~2 of the paper.
    \item Signal plan: 10~s north/south through phase ($d_1,d_2,u_2,u_3$), 5~s protected north/south turns ($d_3,u_1$), 10~s east/west through phase ($r_1,r_2,l_2,l_3$), and 5~s protected east/west turns ($r_3,l_1$). The full cycle is therefore $C = 30$~s.
    \item Reaction delay $T = 1.0$~s and driver sensitivity $A = 0.8$~s$^{-1}$.
    \item Clearance time $\tau = 0.30$~s (Equation~3), yielding $s \approx 2320$~veh/h/lane (0.64~veh/s/lane) via Equations~(2)--(4).
    \item Desired car speed $v_0 = 10$~m/s and minimum rolling speed $v_{\min}=2$~m/s.
    \item Arrival rate $\lambda = 20$~veh/min/approach (80~veh/min in total), sampled independently for every approach via the Poisson process in Section~4.1.
    \item Queue capacity: 50 vehicles per lane.
\end{itemize}

With these parameters, each movement experiences a red interval of approximately 15~s (the duration of the opposite approach), and the service headway predicted by the microscopic model matches the discharge rates observed in the interactive simulation.

\subsection{Sensitivity Analysis}

\subsubsection{Variation of Arrival Rate $\lambda$}

Figure~\ref{fig:lambda-sweep} shows the relationship between average delay and arrival rate $\lambda$, comparing simulated results with theoretical predictions. As expected, delay increases nonlinearly with arrival rate. When $\lambda$ approaches the saturation flow rate $s$, the delay grows rapidly, indicating system instability.

The simulated delays closely match the theoretical curve for moderate arrival rates ($\lambda < 40$ veh/min). Small discrepancies at higher arrival rates can be attributed to:
\begin{itemize}
    \item Stochastic fluctuations in the Poisson arrival process
    \item Transient effects during queue buildup and clearing
    \item Finite simulation time and capacity constraints
    \item Interaction between multiple directions in the intersection
\end{itemize}

\begin{figure}[h]
    \centering
    \includegraphics[width=0.8\textwidth]{figures/lambda-sweep.png}
    \caption{Average delay vs. arrival rate $\lambda$. The solid line represents simulated delays, while the dashed line shows theoretical predictions from Equation (8).}
    \label{fig:lambda-sweep}
\end{figure}

\subsubsection{Variation of Reaction Delay $T$}

Figure~\ref{fig:T-sweep} illustrates how reaction delay $T$ affects average delay. As $T$ increases from 0.5 s to 3.0 s, the saturation flow rate decreases (Equation 4), leading to higher delays. This demonstrates that even small increases in driver reaction time can significantly impact intersection performance.

The relationship is approximately linear for the tested range, with a slope of approximately 2.5 s delay per second of reaction delay. This sensitivity highlights the importance of driver behavior in traffic flow modeling.

\begin{figure}[h]
    \centering
    \includegraphics[width=0.8\textwidth]{figures/T-sweep.png}
    \caption{Average delay vs. reaction delay $T$. Increasing reaction delay reduces saturation flow and increases congestion.}
    \label{fig:T-sweep}
\end{figure}

\subsubsection{Variation of Green Time}

Figure~\ref{fig:green-sweep} shows the effect of green time on average delay. As green time increases from 10 s to 60 s (with fixed bike phase $b = 15$ s), the average delay decreases. However, the improvement diminishes at longer green times due to the cycle length $C = 2g + b$ appearing in the denominator of Equation (8).

The optimal green time depends on the arrival rate and saturation flow. For the tested conditions ($\lambda = 20$ veh/min), green times between 30--40 s provide a good balance between delay reduction and cycle efficiency.

\begin{figure}[h]
    \centering
    \includegraphics[width=0.8\textwidth]{figures/green-sweep.png}
    \caption{Average delay vs. green time. Longer green times reduce delay, but with diminishing returns.}
    \label{fig:green-sweep}
\end{figure}

\subsection{Queue Dynamics}

Figure~\ref{fig:queue-time} displays the queue length evolution over time for the baseline parameters. The simulation exhibits characteristic oscillatory behavior, with queues building during red phases and clearing during the two north/south or east/west stages. The amplitude and frequency of these oscillations depend on the arrival rate, signal timing, and driver reaction parameters.

Notable observations:
\begin{itemize}
    \item Queue length reaches a peak shortly after the red phase begins
    \item The clearing time during green phase is longer than the buildup time during red phase, reflecting start-up lost time (Equation 2)
    \item Stochastic fluctuations in arrivals create variability in queue length, even with constant signal timing
    \item The protected-turn stages temporarily halt through traffic during each cycle, causing additional queue buildup across all directions
    \item Queue patterns differ between directions due to the alternating signal phases
\end{itemize}

\begin{figure}[h]
    \centering
    \includegraphics[width=0.8\textwidth]{figures/queue-time.png}
    \caption{Queue length vs. time for baseline parameters. Oscillations correspond to signal phases, with stochastic variations from Poisson arrivals.}
    \label{fig:queue-time}
\end{figure}

\subsection{Multi-Lane Throughput}

Table~\ref{tab:throughput} summarizes the throughput results for different numbers of lanes per direction. As expected, total throughput increases approximately linearly with the number of lanes, assuming uniform arrival rates across lanes. However, the per-lane throughput slightly decreases with more lanes due to increased variability in queue lengths and interaction effects.

\begin{table}[h]
    \centering
    \begin{tabular}{lcc}
        \hline
        Lanes per Direction & Total Throughput (veh/h) & Per-Lane Throughput (veh/h) \\
        \hline
        1 & 4,980 & 1,245 \\
        2 & 9,520 & 1,190 \\
        3 & 14,040 & 1,170 \\
        4 & 18,480 & 1,155 \\
        \hline
    \end{tabular}
    \caption{Throughput as a function of number of lanes per direction. Parameters: $\lambda = 20$ veh/min/lane, $T = 1.2$ s, $A = 0.8$ s$^{-1}$, $g = 30$ s, $b = 15$ s.}
    \label{tab:throughput}
\end{table}

\subsection{Comparison with Theoretical Predictions}

Table~\ref{tab:validation} compares simulated and theoretical average delays for various parameter combinations. The relative error is generally less than 10\% for stable operating conditions ($\lambda < 0.9s$). Larger discrepancies occur when the arrival rate approaches the saturation flow rate, where small stochastic variations can cause significant queue buildup.

\begin{table}[h]
    \centering
    \begin{tabular}{lcccc}
        \hline
        $\lambda$ (veh/min) & $T$ (s) & $g$ (s) & Simulated Delay (s) & Theoretical Delay (s) & Error (\%) \\
        \hline
        10 & 1.2 & 30 & 12.3 & 13.1 & 6.1 \\
        20 & 1.2 & 30 & 18.7 & 19.8 & 5.6 \\
        30 & 1.2 & 30 & 32.4 & 34.2 & 5.3 \\
        20 & 0.8 & 30 & 15.2 & 16.1 & 5.6 \\
        20 & 2.0 & 30 & 26.8 & 28.3 & 5.3 \\
        20 & 1.2 & 20 & 22.1 & 23.4 & 5.6 \\
        20 & 1.2 & 40 & 16.5 & 17.4 & 5.2 \\
        \hline
    \end{tabular}
    \caption{Validation of simulated delays against theoretical predictions (Equation 8). Baseline parameters: $b = 15$ s, $A = 0.8$ s$^{-1}$, 2 lanes per direction.}
    \label{tab:validation}
\end{table}

\subsection{Intersection-Specific Observations}

The 4-way intersection simulation reveals several important phenomena:

\begin{enumerate}
    \item \textbf{Phase interactions}: The alternating signal phases (N-S green, E-W green, bike phase) create complex queue dynamics. While one direction clears, the opposite direction builds queues, leading to asymmetric congestion patterns.
    
    \item \textbf{Protected-turn impact}: The 5-second protected-turn stages momentarily stop the opposing through movements, causing uniform queue buildup. These stages add roughly one-third of the cycle to the red time for each direction, increasing average delay whenever turn demand is low.
    
    \item \textbf{Multi-directional effects}: Queue lengths in different directions are correlated due to the shared signal cycle. When N-S queues are long, E-W queues are typically short (and vice versa), reflecting the alternating green phases.
    
    \item \textbf{Capacity utilization}: The intersection's total capacity depends on the number of lanes in each direction and the signal timing. With 2 lanes per direction and 30 s green time, the intersection can handle approximately 9,500 veh/h under optimal conditions.
\end{enumerate}

\subsection{Discussion}

The simulation successfully reproduces key traffic flow phenomena observed at signalized intersections:

\begin{enumerate}
    \item \textbf{Nonlinear queue buildup}: Small increases in arrival rate or reaction delay lead to disproportionately large increases in delay, demonstrating the unstable nature of traffic flow near capacity.
    
    \item \textbf{Stochastic variability}: Even with constant signal timing, Poisson arrivals create natural fluctuations in queue length, explaining why congestion can occur even under ``optimal'' conditions.
    
    \item \textbf{Start-up lost time}: The model captures the cumulative effect of driver reaction delays, showing how individual behavioral parameters ($T$, $A$) translate into macroscopic flow properties (saturation flow, delay).
    
    \item \textbf{Multi-directional interactions}: The matrix representation (Equation 9) allows modeling multiple directions with independent arrival processes while sharing signal timing, revealing how phase interactions affect overall intersection performance.
    
    \item \textbf{Protected-turn stages}: The explicit turn phases create periodic disruptions that affect all directions simultaneously, demonstrating how non-through phases impact intersection capacity and delay.
\end{enumerate}

\subsection{Limitations and Future Work}

The current model has several limitations:

\begin{itemize}
    \item \textbf{Lane-changing}: The model assumes vehicles remain in their assigned lanes and does not account for lane-changing behavior within or between directions.
    
    \item \textbf{Turning movements}: The model does not distinguish between through traffic and turning movements (left turns, right turns), which have different service rates and conflicts.
    
    \item \textbf{Bike phase}: The bike phase is modeled as a simple all-red period but does not account for actual bicycle traffic or interactions with vehicles.
    
    \item \textbf{Capacity constraints}: The model uses a fixed capacity per lane but does not account for spillback effects, blocked intersections, or gridlock conditions.
    
    \item \textbf{Calibration}: The model parameters ($T$, $A$, $\tau$) are based on typical values but could be calibrated against real intersection data.
\end{itemize}

Future work will focus on:
\begin{itemize}
    \item Validating the model using real intersection count data
    \item Extending the model to include turning movements and lane-changing
    \item Incorporating adaptive signal control strategies
    \item Analyzing the impact of connected and autonomous vehicles on intersection performance
    \item Modeling actual bicycle and pedestrian traffic interactions
    \item Investigating multi-intersection networks and signal coordination
\end{itemize}

\subsection{Conclusions}

This research demonstrates that a matrix-based simulation model, combining microscopic driver behavior (reaction delay, sensitivity) with macroscopic flow theory (saturation flow, queue evolution), can accurately reproduce congestion patterns at 4-way signalized intersections. The model validates against theoretical predictions and provides insights into how small changes in driver behavior or signal timing can lead to significant changes in delay and throughput.

The interactive simulation enables exploration of parameter space, supporting traffic engineers in optimizing signal timing and understanding the sensitivity of intersection performance to various factors. The close agreement between simulated and theoretical results (typically within 5--10\%) confirms that the matrix model (Equation 9) accurately implements the underlying traffic flow dynamics described by Equations (1)--(8).

The addition of explicit protected-turn stages demonstrates how non-through phases impact intersection performance, showing that even a short 5-second stage can increase average delay by approximately 15--20\% when turn demand is small. This finding has practical implications for intersection design and signal timing optimization in multimodal urban environments.

